
%\needspace{4\baselineskip}
\color{unidarkgreen}{\subsection{Токовая направленная защита нулевой последовательности (ТНЗНП)}}
\color{black}

\begin{enumerate}[label=\arabic{section}.\arabic{subsection}.\arabic{enumi}, labelsep=4pt, leftmargin=0pt, itemindent=57pt, itemsep=0pt, parsep=5pt] % если что можно крутить itemsep

%%%%%%%%%%%%%%%%%%%%%%%%%%%%%%%%%%%%%%%%%%%%%%%%%%%%%%%%%%%%%%%%%%%%%%  ОБЩИЕ СВЕДЕНИЯ %%%%%%%%%%%%%%%%%%%%%%%%%%%%%%%%%%%%%%%%%%%%%%%%%%
\item Общие сведения
\begin{enumerate}[label=\arabic{section}.\arabic{subsection}.\arabic{enumi}.\arabic*, labelsep=4pt, leftmargin=0em, itemindent=65pt, parsep=0pt]

\item
Токовая направленная защита нулевой последовательности (ТНЗНП) предназначена для селективного отключения поврежденного участка электрической сети при коротких замыканиях на землю, реагируя на повышение тока нулевой последовательности.

\item
В устройстве реализовано шесть ступеней ТНЗНП, ввод в работу которых осуществляется программными ключами \textbf{«ТНЗНП-N»}, где N – номер ступени ТНЗНП.

\item
Ток срабатывания ступеней ТНЗНП задается с помощью уставок \textbf{«3I0ср ТНЗНП-N»}. Действие всех ступеней осуществляется с регулируемыми независимыми выдержками времени \textbf{«Тср ТНЗНП-N»}.

\item
На рисунке \ref{pict_tnznp:TNZNP} приведена функциональная схема ТНЗНП-1, выполнение функциональных схем ступеней ТНЗНП-2(3,4,5,6) аналогичное.

\vspace{2mm}
\begin{figure}[H]
\centering
\includegraphics[width=1\textwidth,height=1\textheight,keepaspectratio]{\fbpath/561.0111 ТНЗНП/img/TNZNP.pdf}
\captionof{figure}{Функциональная схема ступени ТНЗНП-1}\label{pict_tnznp:TNZNP}
\end{figure}

\item
Уставки ступеней ТНЗНП представлены в таблице \ref{app_set:tznp}.

\item
Имеется возможность задать различные режимы работы ступеней ТНЗНП при возникновении неисправности в цепях напряжения. Вариант действия логики ступени ТНЗНП при срабатывании БНН задается программным переключателем \textbf{«БНН ТНЗНП-N»}, имеющим три положения:

\begin{enumerate}[label=\arabic*)]
\item \textit{Откл.} -- режим работы ступени не зависит от БНН;
\item \textit{ОНМ} -- ступень переводится в ненаправленный режим работы;
\item \textit{ТНЗНП} --  при срабатывании БНН ступень выводится из работы (однако, если ступень задана ненаправленной, то она останется в работе независимо от срабатывания БНН; ступень также не будет выводиться из работы, если направленность автоматически вывелась при срабатывании ТНЗНП или в процессе автоматического ускорения).
\end{enumerate}

\item
Предусмотрен вывод чувствительных ступеней ТНЗНП, который производится от внешнего сигнала \textit{«ЧС ТНЗНП выведены»}. Программными ключами \textit{«Вывод ТНЗНП-N по чувств.»} можно задать ступени, выводимые из работы.

\item
Предусмотрен автоматический вывод направленности автоматически ускоряемой ступени ТНЗНП (вводится программным ключом \textbf{«Вывод напр. при АУ»}). При наличии разрешающего сигнала ввода ускорения при включении \textit{«Ввод АУ»} направленность ТНЗНП выводится. Причем сигнал вывода направленности подхватывается от пуска АУ ТНЗНП, что обеспечивает устойчивое состояние срабатывания ТНЗНП при неполнофазном включении выключателя. 

\item
Кроме того, направленность может выводиться при срабатывании ТНЗНП. Это обеспечивает устойчивое срабатывание ТНЗНП при неполнофазном отключении, что необходимо для надежного действия УРОВ. Данный режим вводится программным ключом \textbf{«Вывод напр. при ср. ТНЗНП»}.

\item
Для обеспечения несрабатывания ТНЗНП от бросков намагничивающего тока предусмотрена блокировка ступеней ТНЗНП по содержанию второй гармоники в токе нулевой последовательности. 
Блокировка основана на контроле отношения действующего значения второй гармоники тока нулевой последовательности к основной гармонике. Ступень ТНЗНП блокируется при превышении данного соотношения пороговой величины, заданной уставкой \textbf{«Кблок 2 гарм.»}. Наличие функции блокировки от бросков намагничивающего тока задается для каждой ступени отдельно с помощью программного ключа \textbf{«БНТ ТНЗНП-N»}. Для предотвращения излишней блокировки при больших токах КЗ, сопровождающихся насыщением ТТ, блокировка выводится из действия при превышении тока $3I_0$ основной гармоники уставки \textbf{«3I0блок ОВ БНТ»}.

\item
Функциональная схема логики органа выявления БНТ приведена на рисунке \ref{pict_tnznp:bnt}. Уставки органа выявления БНТ представлены в таблице \ref{app_set:tznp}.

\vspace{2mm}
\begin{figure}[H]
\centering
\includegraphics[width=0.8\textwidth,height=0.8\textheight,keepaspectratio]{\fbpath/561.0111 ТНЗНП/img/bnt.pdf}
\captionof{figure}{Функциональная схема логики органа выявления БНТ}\label{pict_tnznp:bnt}
\end{figure}


\item
Для предотвращения неправильной работы защиты предусмотрена блокировка при протекании через трансформаторы тока двух выключателей сквозного тока внешнего КЗ, вызывающего неодновременное насыщение трансформаторов. Блокировка вводится отдельными программными ключами для первой ступени ТНЗНП (\textbf{«БСТО ТНЗНП-1»}), оперативного ускорения (\textbf{«БСТО ОУ ТНЗНП»}), оперативного ускорения без выдержки времени (\textbf{«БСТО ОУ-бв ТНЗНП»}).

\end{enumerate}
%%%%%%%%%%%%%%%%%%%%%%%%%%%%%%%%%%%%%%%%%%%%%%%%%%%%%%%%%%%%%%%%%%%%%%  КОНЕЦ ОБЩИЕ СВЕДЕНИЯ %%%%%%%%%%%%%%%%%%%%%%%%%%%%%%%%%%%%%%%%%%%%%%%%%%

%%%%%%%%%%%%%%%%%%%%%%%%%%%%%%%%%%%%%%%%%%%%%%%%%%%%%%%%%%%%%%%%%%%%%% ОНМ %%%%%%%%%%%%%%%%%%%%%%%%%%%%%%%%%%%%%%%%%%%%%%%%%%
\item Органы направления мощности
\begin{enumerate}[label=\arabic{section}.\arabic{subsection}.\arabic{enumi}.\arabic*, labelsep=4pt, leftmargin=0em, itemindent=65pt, parsep=0pt]

\item
Направленность ТНЗНП обеспечивается органами направления мощности нулевой последовательности (ОНМНП) и обратной последовательности (ОНМОП). 

\item
ОНМНП реагирует на величины векторов тока и напряжения нулевой последовательности, а также угол сдвига между ними. 

\item
Так как мощность нулевой последовательности всегда направлена от места несимметрии во внешнюю сеть, то разрешающий орган направления мощности (ОНМНП-Р) срабатывает при направлении мощности нулевой последовательности от линии к шинам (при КЗ впереди), а блокирующий (ОНМНП-Б) – при обратном направлении мощности. 

\item
За положительное направление отсчета векторов токов и напряжений принято направление против часовой стрелки. Углы максимальной чувствительности откладываются от вектора напряжения $\vec {3U_0}$ (вектор поляризации) и задаются уставками \textbf{«Фмч ОНМНП-Р»} и \textbf{«Фмч ОНМНП-Б»}. Уставки по углам максимальной чувствительности рекомендуется задавать для ОНМНП-Р – плюс $110^\circ$, а для ОНМНП-Б – минус $70^\circ$. Ширина зоны срабатывания $\pm 85^\circ$ от угла максимальной чувствительности.

\item
Пороги срабатывания ОНМНП-Р и ОНМНП-Б по току и напряжению регулируются уставками \textbf{«3I0ср ОНМНП-Р»}, \textbf{«3U0ср ОНМНП-Р»}, \textbf{«3I0ср ОНМНП-Б»}, и \textbf{«3U0ср ОНМНП-Б»}. Выбор напряжения нулевой последовательности, по которому работает ОНМНП, осуществляется с помощью программного переключателя \textbf{«Выбор 3U0»}.

\item
Для повышения чувствительности ОНМНП прямой направленности (ОНМНП-Р) предусмотрена возможность «искусственного смещения точки подключения ТН в линию». Напряжение нулевой 
последовательности, подводимое к ОНМНП-Р, $\vec 3U_{0\textsl{ ОНМНП}}$, В, рассчитывается с учетом добавочного напряжения смещения по формуле:

\begin{equation}
\label{eq:3U}
3\vec U_{0\textsl{ ОНМНП}} = 3\vec U_{0\textsl{ КЗ}} + 3\vec U_{0\textsl{ см}} = 3\vec U_{0\textsl{ КЗ}} + 3 \vec I_{0 \textsl{ КЗ}} \cdot \vec{Z}_{0\textsl{ см}}
\end{equation}
\begin{eqexpl}[10mm]
\item{$3\vec U_{0\textsl{ КЗ}}$}  измеряемое устройством напряжение нулевой последовательности в момент КЗ, во вторичных величинах, В;
\item{$3\vec I_{0 \textsl{ КЗ}}$}  измеряемый устройством ток нулевой последовательности в момент КЗ, во вторичных величинах, А;
\item{$3\vec U_{0\textsl{ см}}$}  напряжение смещения нулевой последовательности, во вторичных величинах, В;
\item{$\vec{Z}_{0\textsl{ см}}$}  сопротивление смещения в линию (вторичное значение в схеме нулевой последовательности), Ом.
\end{eqexpl}

\item
Для смещения точки подключения ТН с шин в линию аргумент сопротивления смещения необходимо принять на $180^\circ$ больше угла максимальной чувствительности. Сопротивление смещения задается уставками \textbf{«Rсм ОНМНП-Р»} и \textbf{«Хсм ОНМНП-Р»}.

\item
ОНМОП реагирует на величины векторов тока и напряжения обратной последовательности, а также угол сдвига между ними.

\item
Разрешающий орган направления мощности обратной последовательности (ОНМОП-Р) срабатывает при направлении мощности обратной последовательности от линии к шинам, а блокирующий (ОНМОП-Б) -- при обратном направлении мощности. 

\item
Углы максимальной чувствительности задаются уставками \textbf{«Фмч ОНМОП-Р»} и \textbf{«Фмч ОНМОП-Б»}. Уставки по углам максимальной чувствительности рекомендуется задавать для ОНМОП-Р -- плюс $110^\circ$, а для ОНМОП-Б -- минус $70^\circ$. Ширина зоны срабатывания $\pm 85^\circ$ от угла максимальной чувствительности.

\item
Пороги срабатывания ОНМОП-Р и ОНМОП-Б по току и напряжению регулируются уставками \textbf{«3I0ср ОНМОП-Р»}, \textbf{«3U0ср ОНМОП-Р»}, \textbf{«3I0ср ОНМОП-Б»}, и \textbf{«3U0ср ОНМОП-Б»}.

\item
Параметры ОНМНП(ОНМОП) приведены в таблице \ref{desc_tnznp:tbl2}.

\small
\setlength{\extrarowheight}{0.05cm} % высота ячеек в таблице
\begin{longtable}
{|>{\centering\arraybackslash}m{15.2cm}|>{\centering\arraybackslash}m{1.7cm}|}
\caption{Параметры ОНМНП(ОНМОП)\hfill\vspace{-0.5\baselineskip}}\label{desc_tnznp:tbl2}\\ 
\hhline{|-|-|}
\rowcolor{gray!30}
Наименование параметра & Значение  \\ 
\hhline{|-|-|}
\endfirsthead

\caption*{\hspace{3pt}\emph{Продолжение таблицы \ref{desc_tnznp:tbl2}\hfill\vspace{-0.5\baselineskip}}} \\ % сделано по ГОСТ 2.105 п.6.8.7
\hhline{|-|-|}
\rowcolor{gray!30}
Наименование параметра & Значение  \\ 
\hhline{|-|-|}
\endhead
\hhline{|-|-|}

\endfoot
\endlastfoot

\raggedright Время срабатывания ОНМ при утроенных по отношению к порогам срабатывания значениях тока и напряжения, мс, не более & \centering\arraybackslash 35  \\  \hline
\raggedright Время возврата ОНМ при одновременном сбросе входных величин тока и напряжения от номинальных значений до нуля, мс, не более  & \centering\arraybackslash 25  \\  \hline
\raggedright Средняя основная погрешность порогов срабатывания ОНМ по току и напряжению, \% уставки, не более   & \centering\arraybackslash $\pm$5  \\  \hline
\raggedright Средняя основная абсолютная погрешность ОНМ по углу максимальной чувствительности, эл. град., не более & \centering\arraybackslash $\pm$5 \\  \hline
\raggedright Дополнительная погрешность параметров срабатывания от изменения температуры окружающего воздуха в рабочем диапазоне, \% от среднего значения параметров, измеренных при температуре $(20\pm5)^\circ$C, не более & \centering\arraybackslash $\pm$5 \\  \hline
\raggedright Минимальная угловая ширина зоны срабатывания ОНМ, эл. град, не менее & \centering\arraybackslash 165 \\ \hline
\raggedright Коэффициент возврата ОНМ НП по напряжению и току, не менее  & \centering\arraybackslash 0,9\\  \hline

\end{longtable}
\setlength{\extrarowheight}{0.0cm} % возврат параметра высота ячеек в таблице
\normalsize

\item
На рисунке \ref{pict_tnznp:char_triggering_ONMNP} представлены характеристики срабатывания ОНМНП-Р и ОНМНП-Б. Аналогичные характеристики для ОНМОП-Р и ОНМОП-Б.

\vspace{2mm}
\begin{figure}[H]
\centering
\includegraphics[width=1\textwidth,height=1\textheight,keepaspectratio]{\fbpath/561.0111 ТНЗНП/img/char_triggering_ONMNP.pdf}
\captionof{figure}{Характеристики срабатывания ОНМНП}\label{pict_tnznp:char_triggering_ONMNP}
\end{figure}

\item
Для выбора режима направленности по мощности нулевой последовательности для каждой ступени ТНЗНП предусмотрен программный переключатель \textbf{«Направл. ТНЗНП-N»}, имеющий три положения:

\begin{enumerate}[label=\arabic*)]
    \item \textit{Ненапр.} -- ступень выполнена ненаправленной;
    \item \textit{Прямо} -- ступень выполнена прямонаправленной;
    \item \textit{Обратно} -- ступень выполнена обратнонаправленной.
\end{enumerate}

\item
Кроме того, с помощью программного переключателя \textbf{«Контр. направл. ТНЗНП-N»} можно установить различные алгоритмы контроля направленности.

\begin{enumerate}[label=\arabic*)]
    \item \textit{Разр.} --  прямая направленность обеспечивается срабатыванием ОНМНП-Р, обратная -- срабатыванием ОНМНП-Б;
    \item \textit{Блок.} -- прямая направленность обеспечивается отсутствием срабатывания ОНМНП-Б, обратная -- отсутствием срабатывания ОНМНП-Р;
    \item \textit{Разр. или блок.} -- прямая направленность обеспечивается срабатыванием ОНМНП-Р или отсутствием срабатывания ОНМНП-Б, обратная -- срабатыванием ОНМНП-Б или отсутствием срабатывания ОНМНП-Р.
\end{enumerate}

\item
Положение переключателя \textit{«Разр. или блок.»} рекомендуется выставлять для ступеней ТНЗНП, выполняющих функцию дальнего резервирования, так как при удаленных КЗ напряжение $3U_0$ может оказаться недостаточным для срабатывания ОНМ из-за его уменьшения по мере удаления от точки КЗ. Уставки для выбора режима направленности по мощности обратной последовательности аналогичные.

\item
Тип ОНМ, используемый для определения направленности ступени ТНЗНП, выбирается программным переключателем \textit{«Тип ОНМ ТНЗНП-N»}, имеющим четыре положения:

\begin{enumerate}[label=\arabic*)]
    \item \textit{НП} -- для определения направленности ступени учитывается только ОНМНП;
    \item \textit{ОП} -- для определения направленности ступени учитывается только ОНМОП;
    \item \textit{НП или ОП} -- для определения направленности ступени учитывается один из органов, ОНМНП или ОНМОП;
    \item \textit{НП и ОП} -- для определения направленности ступени одновременно учитываются оба органа, ОНМНП и ОНМОП.
\end{enumerate}

\item
Функциональная схема логики ОНМНП приведена на рисунке \ref{pict_tnznp:onnp}. Функциональная схема логики ОНМОП приведена на рисунке \ref{pict_tnznp:onop}. Уставки ОНМОП и ОНМНП представлены в таблице \ref{app_set:tznp}.

\vspace{2mm}
\begin{figure}[H]
\centering
\includegraphics[width=0.8\textwidth,height=0.8\textheight,keepaspectratio]{\fbpath/561.0111 ТНЗНП/img/onnp.pdf}
\captionof{figure}{Функциональная схема логики ОНМНП}\label{pict_tnznp:onnp}
\end{figure}

\vspace{2mm}
\begin{figure}[H]
\centering
\includegraphics[width=0.8\textwidth,height=0.8\textheight,keepaspectratio]{\fbpath/561.0111 ТНЗНП/img/onop.pdf}
\captionof{figure}{Функциональная схема логики ОНМОП}\label{pict_tnznp:onop}
\end{figure}


\end{enumerate}

%%%%%%%%%%%%%%%%%%%%%%%%%%%%%%%%%%%%%%%%%%%%%%%%%%%%%%%%%%%%%%%%%%%%%% АУ %%%%%%%%%%%%%%%%%%%%%%%%%%%%%%%%%%%%%%%%%%%%%%%%%%
\item Автоматическое ускорение (АУ)
\begin{enumerate}[label=\arabic{section}.\arabic{subsection}.\arabic{enumi}.\arabic*, labelsep=4pt, leftmargin=0em, itemindent=65pt, parsep=0pt]

\item
В составе ТНЗНП предусмотрена логика формирования пуска автоматического ускорения ТНЗНП (АУ ТНЗНП), выходной сигнал которого действует в логику автоматического ускорения. Формирование сигнала пуска АУ ТНЗНП можно задать от любой ступени ТНЗНП программным переключателем \textbf{«Ступень АУ ТНЗНП»}. Сигнал пуска АУ ТНЗНП формируется непосредственно от сигналов пуска ступени ТНЗНП, то есть настройка автоматически ускоряемой ступени определяется соответствующими настройками самой ступени.

\item
Функциональная схема логики АУ ТНЗНП приведена на рисунке \ref{pict_tnznp:AU_TNZNP}.

\vspace{2mm}
\begin{figure}[H]
\centering
\includegraphics[width=1\textwidth,height=1\textheight,keepaspectratio]{\fbpath/561.0111 ТНЗНП/img/AU_TNZNP.pdf}
\captionof{figure}{Функциональная схема пуска АУ ТНЗНП}\label{pict_tnznp:AU_TNZNP}
\end{figure}

\item
Уставки АУ ТНЗНП представлены в таблице \ref{app_set:tznp}.


\end{enumerate}

%%%%%%%%%%%%%%%%%%%%%%%%%%%%%%%%%%%%%%%%%%%%%%%%%%%%%%%%%%%%%%%%%%%%%% ОУ %%%%%%%%%%%%%%%%%%%%%%%%%%%%%%%%%%%%%%%%%%%%%%%%%%
\item Оперативное ускорение (ОУ)
\begin{enumerate}[label=\arabic{section}.\arabic{subsection}.\arabic{enumi}.\arabic*, labelsep=4pt, leftmargin=0em, itemindent=65pt, parsep=0pt]

\item
В устройстве реализовано два режима оперативного ускорения ТНЗНП: с выдержкой времени (ОУ ТНЗНП) и без выдержки времени (ОУ-бв ТНЗНП). Ввод в работу ОУ ТНЗНП осуществляется от внешнего сигнала \textit{«Ввод ОУ ТНЗНП»}. Выбор оперативно ускоряемой ступени ТНЗНП выполняется с помощью программного переключателя \textbf{«Ступень ОУ ТНЗНП»}. Время срабатывания ОУ ТНЗНП задается уставкой \textbf{«Тср ОУ ТНЗНП»}.

\item
Функциональная схема логики ОУ ТНЗНП приведена на рисунке \ref{pict_tnznp:OU_TNZNP}. Функциональная схема логики ОУ-бв ТНЗНП выполнена аналогично ОУ ТНЗНП.

\vspace{2mm}
\begin{figure}[H]
\centering
\includegraphics[width=1\textwidth,height=1\textheight,keepaspectratio]{\fbpath/561.0111 ТНЗНП/img/OU_TNZNP.pdf}
\captionof{figure}{Функциональная схема ОУ ТНЗНП}\label{pict_tnznp:OU_TNZNP}
\end{figure}

\item
Ввод в работу ОУ-бв ТНЗНП осуществляется от внешнего сигнала \textit{«Ввод ОУ-бв ТНЗНП»}. Выбор оперативно ускоряемой ступени ТНЗНП без выдержки времени выполняется с помощью программного переключателя \textbf{«Ступень ОУ-бв ТНЗНП»}. При необходимости для ОУ-бв ТНЗНП также можно задать время срабатывания уставкой \textbf{«Тср ОУ-бв ТНЗНП»}.

\item
Кроме того, предусмотрена возможность автоматического ввода ОУ ТНЗНП и ОУ-бв ТНЗНП программными ключами \textbf{«Ввод ОУ ТНЗНП при выводе осн. защиты»} и \textbf{«Ввод ОУ-бв ТНЗНП при выводе осн. защиты»} при приеме внешнего сигнала о выводе основной защиты. Введенное ускорение может быть оперативно выведено от внешнего сигнала \textit{«Авт. введенное ОУ выведено»}.

\item
Уставки ОУ ТНЗНП и оперативного ускорения ТНЗНП без выдержки времени (ОУ-бв ТНЗНП) представлены в таблице \ref{app_set:tznp}.

\end{enumerate}

%%%%%%%%%%%%%%%%%%%%%%%%%%%%%%%%%%%%%%%%%%%%%%%%%%%%%%%%%%%%%%%%%%%%%% ПУ %%%%%%%%%%%%%%%%%%%%%%%%%%%%%%%%%%%%%%%%%%%%%%%%%%
\item Поперечное ускорение (ПУ)
\begin{enumerate}[label=\arabic{section}.\arabic{subsection}.\arabic{enumi}.\arabic*, labelsep=4pt, leftmargin=0em, itemindent=65pt, parsep=0pt]

\item
В устройстве предусмотрено поперечное ускорение ТНЗНП (ПУ ТНЗНП) от защит параллельной линии, принцип действия которого основан на сравнении направления мощности нулевой последовательности защищаемой и параллельной линии.

\item
Программный переключатель \textbf{«Ступень ПУ ТНЗНП»} задает ускоряемую ступень ТНЗНП.

\item
Функциональная схема логики ПУ ТНЗНП приведена на рисунке \ref{pict_tnznp:PU_TNZNP}.

\vspace{2mm}
\begin{figure}[H]
\centering
\includegraphics[width=1\textwidth,height=1\textheight,keepaspectratio]{\fbpath/561.0111 ТНЗНП/img/PU_TNZNP.pdf}
\captionof{figure}{Функциональная схема пуска ПУ ТНЗНП}\label{pict_tnznpp:PU_TNZNP}
\end{figure}

\item
Уставки ПУ ТНЗНП представлены в таблице \ref{app_set:tznp}.

\item
Данный вид ускорения может корректно работать только при наличии связи между параллельными линиями. Если шины соединены через ШСВ, то необходим контроль его включенного состояния (сигнал \textit{«ШСВ/СВ включен»}). Если ШСВ выведен в ремонт или отсутствует, но имеется связь между параллельными линиями с помощью оперативного переключателя подается дискретный сигнал «ШСВ/СВ выведен».

\item
В таблице \ref{desc_tnznp:tbl8} приведен пример использования внешнего оперативного переключателя для задания режима работы ускорения в соответствии с состоянием ШСВ.

\small
\setlength{\extrarowheight}{0.05cm} % высота ячеек в таблице
\begin{longtable}
{|>{\centering\arraybackslash}m{4cm}|>{\centering\arraybackslash}m{4cm}|>{\centering\arraybackslash}m{4cm}|>{\centering\arraybackslash}m{4cm}|}
\caption{Формируемые сигналы в устройстве в зависимости от положения переключателя:\hfill\vspace{-0.5\baselineskip}}\label{desc_tnznp:tbl8}\\ 
\hhline{|-|-|-|-|}
\rowcolor{gray!30}
№ положения & Наименование  & Вход ФСУ \textit{«ШСВ/СВ выведен»} & Вход ФСУ \textit{«Вывод ПУ ТНЗНП»}  \\ 
\hhline{|-|-|-|-|}
\endfirsthead

\caption*{\hspace{3pt}\emph{Продолжение таблицы \ref{desc_tnznp:tbl8}\hfill\vspace{-0.5\baselineskip}}} \\ % сделано по ГОСТ 2.105 п.6.8.7
\hhline{|-|-|-|-|}
\rowcolor{gray!30}
№ положения & Наименование & Вход ФСУ \textit{«ШСВ/СВ выведен»} & Вход ФСУ \textit{«Вывод ПУ ТНЗНП»} \\ 
\hhline{|-|-|-|-|}
\endhead
\hhline{|-|-|-|-|}

\endfoot
\endlastfoot

\raggedright \centering\arraybackslash 1 & ШСВ выведен & 1 & 0 \\  \hline
\raggedright \centering\arraybackslash 2 & Выведено & 0 & 1 \\  \hline
\raggedright \centering\arraybackslash 3 & ШСВ в работе & 0 & 0 \\  \hline

\end{longtable}
\setlength{\extrarowheight}{0.0cm} % возврат параметра высота ячеек в таблице
\normalsize

\item
Если ШСВ выведен, то для формирования сигнала срабатывания ПУ ТНЗНП достаточно приёма сигнала \textit{«ПУ от пар. ЛЭП»} (сигнал срабатывания ОНМНП-Б и РПВ параллельной линии), что свидетельствует об отсутствии повреждения на параллельной линии. Использование сигнала РПВ параллельной ВЛ позволяет исключить неправильное действие защиты при повреждении на параллельной ВЛ в зоне между выносными трансформаторами тока и одним из выключателей этой линии.

\item
Если переключатель в положении \textit{«Выведено»}, то логика ПУ ТНЗНП выводится.
\item
Если ШСВ в работе, то необходимо контролировать его включенное положение, в данном случае для 
формирования сигнала срабатывания ПУ ТНЗНП необходимо наличие двух сигналов: \textit{«ПУ от пар. ЛЭП»} и 
\textit{«ШСВ/CВ включен»}.

\item
Время срабатывания поперечного ускорения ТНЗНП определяется уставкой \textbf{«Тср ПУ ТНЗНП»}.

\item
Формирование сигнала \textit{«ПУ ТНЗНП пар. ЛЭП»} предназначено для выдачи на устройство РЗА, установленное на параллельной ЛЭП. Он то и должен назначаться на вход ФСУ \textit{«ПУ от пар. ЛЭП»} устройства на параллельной линии. 

\end{enumerate}

\end{enumerate}